% LaTeX document recreating the French learning guide
\documentclass[a4paper,12pt]{article}
\usepackage[utf8]{inputenc}
\usepackage[T1]{fontenc}
\usepackage[brazil]{babel}
\usepackage{geometry}
\geometry{a4paper, margin=2.5cm}
\usepackage{array}
\usepackage{tabularx}
\usepackage{booktabs}
\usepackage{enumitem}
\usepackage{xcolor}
\usepackage{hyperref}
\hypersetup{
    colorlinks=true,
    linkcolor=blue,
    urlcolor=blue,
}
\usepackage{bookmark}  % added to avoid hyperref warnings
\usepackage{graphicx}
\usepackage{tikz}
\usepackage{longtable}
\usepackage{multicol}

% For conditional formatting and better tables
\newcolumntype{L}{>{\raggedright\arraybackslash}X}

% Title and author
\title{Francês --- Guia Completo de Aprendizagem}
\author{}
\date{Última atualização: 6 de março de 2026}

\begin{document}

\maketitle

\section*{Objetivo}
B1 (Conversação prática) em 200--240 horas \\
\textbf{Timeline:} 4 meses com 45--90 min/dia, 5--6 dias/semana \\
\textbf{Método Base:} Active Recall (Anki) + Aulas estruturadas + Conversação \\
\textbf{Evidência Científica:} Spaced Repetition (Ebbinghaus 1880) + Active Recall (50+ estudos)

\section*{Nota importante}
Este guia é estruturado com base em pesquisas de aquisição de linguagem revisadas por pares:
\begin{itemize}
    \item CEFR Framework: A1 (Iniciante) $\to$ A2 (Elementar) $\to$ B1 (Conversação prática)
    \item Spaced Repetition: Reveja a gramática em intervalos crescentes (1 dia, 3 dias, 7 dias, 14 dias)
    \item Active Recall: Teste-se nas regras gramaticais — NÃO leitura passiva
    \item Sono e Memória: Estude à noite $\to$ 7--9h de sono $\to$ pratique no dia seguinte = 2x taxa de retenção
\end{itemize}

\tableofcontents
\newpage

% ----------------------------------------------------------------------
% Seção: Roadmap Visual
% ----------------------------------------------------------------------
\section{Roadmap Visual — Caminho para B1}

\begin{center}
\begin{tikzpicture}[scale=0.7, every node/.style={align=center}]
    % Draw horizontal timeline
    \draw[thick,->] (0,0) -- (16,0) node[anchor=west] {Semanas};
    \foreach \x in {0,4,8,12,16}
        \draw (\x,0.1) -- (\x,-0.1) node[below] {\x};
    % Level boxes
    \fill[blue!20] (0,0.2) rectangle (4,1) node[midway] {A1\\Fundações};
    \fill[green!20] (4,0.2) rectangle (8,1) node[midway] {A2\\Passado};
    \fill[orange!20] (8,0.2) rectangle (12,1) node[midway] {B1 Inicial\\Futuro};
    \fill[red!20] (12,0.2) rectangle (16,1) node[midway] {B1 Fluência\\Consolidação};
    \draw[thick,<->] (0,1.2) -- (4,1.2) node[midway,above] {40h};
    \draw[thick,<->] (4,1.2) -- (8,1.2) node[midway,above] {80h};
    \draw[thick,<->] (8,1.2) -- (12,1.2) node[midway,above] {130h};
    \draw[thick,<->] (12,1.2) -- (16,1.2) node[midway,above] {200h};
\end{tikzpicture}
\end{center}

\subsection*{Distribuição Total}
\begin{tabular}{clcc}
    \toprule
    Mês & Horas & Nível & Foco Principal \\
    \midrule
    Mês 1 & 40h & A1 & Present regular \& irregular \\
    Mês 2 & 80h & A2 & Passado completo \\
    Mês 3 & 130h & B1 (Inicial) & Futuro \& Conversação \\
    Mês 4 & 200h & B1 & Fluência \\
    \bottomrule
\end{tabular}

% ----------------------------------------------------------------------
% Seção: Progressão de Horas
% ----------------------------------------------------------------------
\section{Progressão de Horas Investidas}

\subsection*{Cronograma Semanal}
\textbf{Distribuição de Tempo:}
\begin{itemize}
    \item Total Base: 260 minutos/semana = 4,3 horas/semana
    \item 16 Semanas: 69 horas de estudo dirigido + 130 horas de exposição = 200 horas totais
    \item Semana 1: 104 min (40\% — ajuste)
    \item Semanas 2--4: 260 min (60\% — consistente)
    \item Semanas 5--8: 260 min (60\% — consolidação)
    \item Semanas 9--16: 280 min (70\% — foco intenso)
\end{itemize}

\subsection*{Progresso Cumulativo por Mês}
\begin{tabular}{clcc}
    \toprule
    Período & Horas Acumuladas & Nível & Domínio \\
    \midrule
    Mês 1 & 40h & A1 & Present (regular \& irregular) \\
    Mês 2 & 80h & A2 & Passado completo \\
    Mês 3 & 130h & B1 Inicial & Futuro + Conversação \\
    Mês 4 & 200h & B1 & Fluência \\
    \bottomrule
\end{tabular}
\emph{META: 200--240h}

\subsection*{Timeline de 16 Semanas — Sua Jornada até B1}
\begin{tabular}{clccc}
    \toprule
    Período & Semanas & Nível & Foco & Meta de Horas \\
    \midrule
    Fundações & 1--4 & A1 & Verbos -ER/-IR, Présent Regular & 40h \\
    Passado & 5--8 & A2 & Passé Composé, Imparfait, DRMRS & 80h \\
    Futuro \& Conversação & 9--12 & B1 Inicial & Futur, Conditionnel, Subjonctif & 130h \\
    Consolidação & 13--16 & B1 Fluência & Reading, Listening, Production & 200h \\
    \bottomrule
\end{tabular}

\subsection*{Marcos Importantes}
\begin{itemize}
    \item \textbf{Semana 4:} Domina todos os tempos no presente $\to$ Primeiro teste de conversação simples
    \item \textbf{Semana 8:} Completa o passado $\to$ Pode narrar histórias em francês
    \item \textbf{Semana 12:} Domina futuro e condicional $\to$ Conversa sobre planos e hipóteses
    \item \textbf{Semana 16:} Chega a B1 $\to$ Simula prova DELF com sucesso
\end{itemize}

% ----------------------------------------------------------------------
% Seção: Estrutura Gramatical
% ----------------------------------------------------------------------
\section{Estrutura Gramatical Francesa — Visão Geral}

\subsection{Partes do Discurso (Parts of Speech)}
\subsubsection{Substantivos (Noms)}
\begin{itemize}
    \item Todos os substantivos têm gênero: masculino ou feminino
    \item São precedidos por artigos (\emph{un, une, le, la, les})
    \item Exemplo: \emph{le livre} (o livro - masculino), \emph{la maison} (a casa - feminino)
\end{itemize}

\subsubsection{Artigos (Articles)}
\begin{tabular}{clcc}
    \toprule
    Tipo & Singular & Plural \\
    \midrule
    Indefinido & un (m), une (f) & des \\
    Definido & le (m), la (f), l' (vocal) & les \\
    Partitivo & du, de la, de l' & des \\
    \bottomrule
\end{tabular}

\subsubsection{Adjetivos (Adjectifs)}
\begin{itemize}
    \item Concordam em gênero e número com o substantivo
    \item Feminino geralmente adiciona \emph{-e} ao masculino
    \item Plural geralmente adiciona \emph{-s} ao singular
    \item Exemplo: \emph{un livre rouge} (um livro vermelho), \emph{une maison rouge} (uma casa vermelha)
\end{itemize}

\subsubsection{Verbos (Verbes)}
Três grupos principais baseados na terminação do infinitivo:
\begin{itemize}
    \item 1º grupo: verbos \emph{-ER} (parler, danser) — cerca de 80\% dos verbos
    \item 2º grupo: verbos \emph{-IR} (finir, choisir) — cerca de 300 verbos
    \item 3º grupo: verbos irregulares (être, avoir, aller, venir) — cerca de 100 verbos
\end{itemize}

\subsubsection{Pronomes (Pronoms)}
\begin{tabular}{lcc}
    \toprule
    Tipo & Singular & Plural \\
    \midrule
    Sujeito & je, tu, il, elle, on & nous, vous, ils, elles \\
    Objeto direto & me, te, le, la & nous, vous, les \\
    Objeto indireto & me, te, lui & nous, vous, leur \\
    Reflexivo & me, te, se & nous, vous, se \\
    \bottomrule
\end{tabular}

\subsection{Tempos Verbais Essenciais}
\subsubsection{PRÉSENT (Presente)}
Uso: Ações atuais e habituais. \\
Estrutura Regular -ER: \emph{je parle, tu parles, il/elle parle, nous parlons, vous parlez, ils/elles parlent} \\
Exemplos: \emph{Je parle français} | \emph{Tu danses bien} | \emph{Il chante}

\subsubsection{PASSÉ COMPOSÉ (Perfeito)}
Uso: Ações passadas concluídas. \\
Estrutura: Auxiliar (avoir/être) + Participio Passado. \\
Exemplos:
\begin{itemize}
    \item Com avoir: \emph{J’ai parlé} (Eu falei)
    \item Com être: \emph{Je suis allé(e)} (Eu fui)
\end{itemize}

\subsubsection{IMPARFAIT (Pretérito Imperfeito)}
Uso: Ações passadas contínuas ou habituais; narração de contexto. \\
Estrutura: radical do presente + terminações (-ais, -ais, -ait, -ions, -iez, -aient). \\
Exemplos: \emph{Je parlais} (Eu falava / estava falando) | \emph{Tu dansais} | \emph{Il chantait}

\subsubsection{FUTUR SIMPLE (Futuro Simples)}
Uso: Ações futuras. \\
Estrutura Regular -ER: \emph{je parlerai, tu parleras, il/elle parlera, nous parlerons, vous parlerez, ils/elles parleront} \\
Exemplos: \emph{Je parlerai français} | \emph{Tu danseras demain}

\subsubsection{CONDITIONNEL (Condicional)}
Uso: Situações hipotéticas, polidez, conselhos. \\
Estrutura: Radical do futuro + terminações do imparfait. \\
Exemplos: \emph{Je parlerais} (Eu falaria) | \emph{Je voudrais} (Eu gostaria) | \emph{Tu devrais} (Tu devias)

\subsubsection{SUBJONCTIF (Subjuntivo)}
Uso: Após expressões de dúvida, desejo, necessidade, emoção. \\
Estrutura: que + sujeito + forma subjuntiva. \\
Exemplos: \emph{Je veux que tu parles} (Eu quero que tu fales) | \emph{Il faut que je parte} (É necessário que eu parta)

\subsection{Estrutura da Frase}
\subsubsection{Ordem das Palavras}
Padrão Básico: Sujeito + Verbo + Objeto. \\
Exemplos:
\begin{itemize}
    \item \emph{Je parle français} (Eu falo francês)
    \item \emph{Elle lit un livre} (Ela lê um livro)
    \item \emph{Nous mangeons une pizza} (Nós comemos uma pizza)
\end{itemize}
Perguntas: Inversão (verbo antes do sujeito) ou \emph{est-ce que}.
\begin{itemize}
    \item \emph{Tu parles français?} (Você fala francês? — somente com entonação)
    \item \emph{Parles-tu français?} (Você fala francês? — inversão formal)
    \item \emph{Est-ce que tu parles français?} (Você fala francês? — com est-ce que)
\end{itemize}

\subsubsection{Negação}
\begin{tabular}{lll}
    \toprule
    Estrutura & Significado & Exemplo \\
    \midrule
    ne...pas & não & \emph{Je ne suis pas} (Eu não sou) \\
    ne...jamais & nunca & \emph{Je ne vais jamais} (Eu nunca vou) \\
    ne...rien & nada & \emph{Je ne dis rien} (Eu não digo nada) \\
    ne...plus & não mais & \emph{Je ne suis plus} (Eu já não sou) \\
    \bottomrule
\end{tabular}

\subsubsection{Palavras Interrogativas}
Essenciais para fazer perguntas:
\begin{itemize}
    \item \emph{Qui?} (Quem?) | \emph{Qui est là?}
    \item \emph{Que/Quoi?} (O que?) | \emph{Que fais-tu?}
    \item \emph{Où?} (Onde?) | \emph{Où vas-tu?}
    \item \emph{Quand?} (Quando?) | \emph{Quand arrives-tu?}
    \item \emph{Pourquoi?} (Por quê?) | \emph{Pourquoi souris-tu?}
    \item \emph{Comment?} (Como?) | \emph{Comment vas-tu?}
\end{itemize}

\subsection{Concordância de Gênero e Número}
\begin{itemize}
    \item Gênero: Masculino (\emph{le, un}) / Feminino (\emph{la, une})
    \item Adjetivos e artigos devem concordar com o gênero do substantivo
    \item Número: Singular (um item) / Plural (múltiplos itens, geralmente adiciona \emph{-s})
    \item Plurais irregulares em \emph{-au, -eau, -eux}
\end{itemize}

\subsection{Preposições (Prépositions)}
\begin{tabular}{ll}
    \toprule
    Preposição & Uso \\
    \midrule
    à & a, em \\
    de & de, desde \\
    en & em \\
    sur & sobre \\
    sous & sob \\
    avec & com \\
    sans & sem \\
    entre & entre \\
    devant & diante de \\
    \bottomrule
\end{tabular}
Contrações: \emph{au} (à + le), \emph{du} (de + le), \emph{aux} (à + les), \emph{des} (de + les).

% ----------------------------------------------------------------------
% Seção: Plano de Aulas Detalhado
% ----------------------------------------------------------------------
\section{Plano de Aulas Detalhado Semana a Semana (16 Semanas)}
\subsection{Mês 1: Présent (Semanas 1–4)}
\subsubsection*{Semana 1: Présent dos Verbos -ER (Regulares)}
\begin{tabularx}{\linewidth}{>{\bfseries}l L}
    Gramática & Presente do indicativo: verbos terminados em \emph{-ER} (parler, chanter, danser, étudier, travailler). Conjugação: -e, -es, -e, -ons, -ez, -ent. \\
    Vocabulário & 30 cognatos diretos (palavras idênticas ou muito próximas): animal, important, possible, problème, général, naturel, original, particulier, visible, horrible, terrible, flexible, responsable, art, part, classe, phrase, image, ligne, musique, note, texte, total, triple, universal, vertical, etc. \\
    Exercícios & 1. Conjugar 10 verbos -ER no presente. 2. Traduzir frases simples: "Eu falo francês" $\to$ "\emph{Je parle français}". 3. Identificar cognatos em um texto curto. \\
    Conversação & Apresentar-se: "\emph{Je m'appelle...}", "\emph{Je suis portugais...}", "\emph{Je parle portugais et j'apprends le français}". \\
    Anki & 20 cards: verbos -ER regulares + cognatos. \\
\end{tabularx}

\subsubsection*{Semana 2: Présent dos Verbos -IR e -RE (Regulares)}
\begin{tabularx}{\linewidth}{>{\bfseries}l L}
    Gramática & Presente: verbos -IR (finir, choisir, réussir) — conjugação: -is, -is, -it, -issons, -issez, -issent; verbos -RE (vendre, attendre, entendre) — conjugação: -s, -s, (nada), -ons, -ez, -ent. \\
    Vocabulário & 30 palavras com padrões de conversão: -tion (nation, chanson, éducation), -té (cité, vérité, université), -ble (aimable, possible, terrible). \\
    Exercícios & 1. Conjugar verbos -IR/-RE. 2. Converter palavras portuguesas usando padrões. 3. Completar frases. \\
    Conversação & Falar sobre rotina: "\emph{Tous les jours, je mange à midi...}", "\emph{J'habite à...}", "\emph{Je travaille comme...}" \\
    Anki & 20 cards: verbos -IR/-RE + palavras com padrões. \\
\end{tabularx}

\subsubsection*{Semana 3: Présent dos Verbos Irregulares}
\begin{tabularx}{\linewidth}{>{\bfseries}l L}
    Gramática & Presente dos irregulares fundamentais: \\
    & - AVOIR (j'ai, tu as, il a, nous avons, vous avez, ils ont) \\
    & - ÊTRE (je suis, tu es, il est, nous sommes, vous êtes, ils sont) \\
    & - ALLER (je vais, tu vas, il va, nous allons, vous allez, ils vont) \\
    & - FAIRE (je fais, tu fais, il fait, nous faisons, vous faites, ils font) \\
    Vocabulário & Expressões comuns com esses verbos: avoir faim/soif/froid/chaud, être fatigué/malade, aller bien/mal, faire attention/chaud/froid. \\
    Exercícios & 1. Completar frases com o verbo correto. 2. Traduzir expressões do português. 3. Formar frases negativas. \\
    Conversação & Descrever estados e ações: "\emph{J'ai faim. Je suis fatigué. Je vais au travail. Je fais du sport.}" \\
    Anki & 20 cards: verbos irregulares + expressões. \\
\end{tabularx}

\subsubsection*{Semana 4: Verbos Pronominais + Falsos Amigos + Revisão}
\begin{tabularx}{\linewidth}{>{\bfseries}l L}
    Gramática & Verbos pronominais (reflexivos): se lever, s'appeler, se laver, se promener. Conjugação com pronomes reflexivos: je me lève, tu te lèves, il se lève, nous nous levons, vous vous levez, ils se lèvent. \\
    Vocabulário & Lista dos 30 falsos amigos mais comuns (ver seção específica). Ênfase: actuellement, librairie, présenter, etc. \\
    Exercícios & 1. Conjugar verbos pronominais. 2. Tradução evitando falsos amigos. 3. Pequeno texto usando presente e pronominais. \\
    Conversação & Rotina diária com pronominais: "\emph{Je me lève à 7h, je m'appelle...}" \\
    Anki & 20 cards: falsos amigos + verbos pronominais. \\
\end{tabularx}

\subsection{Mês 2: Passado (Semanas 5–8)}
\subsubsection*{Semana 5: Passé Composé com Avoir (Regulares)}
\begin{tabularx}{\linewidth}{>{\bfseries}l L}
    Gramática & Passé Composé: avoir conjugado no presente + participio passado. Participios regulares: -ER $\to$ -é (parlé), -IR $\to$ -i (fini), -RE $\to$ -u (vendu). Negação: n'ai pas parlé. \\
    Vocabulário & Marcadores temporais: hier, avant-hier, la semaine dernière, en 2020, déjà, jamais. \\
    Exercícios & 1. Formar participios passados. 2. Conjugar verbos no passé composé. 3. Narrar ações concluídas: "\emph{Hier, j'ai mangé une pizza.}" \\
    Conversação & Contar o que fez ontem: "\emph{Hier, je me suis levé à 8h, j'ai travaillé, j'ai déjeuné avec des amis...}" \\
    Anki & 20 cards: passé composé regular. \\
\end{tabularx}

\subsubsection*{Semana 6: Passé Composé com Être (DR \& MRS VANDERTRAMP)}
\begin{tabularx}{\linewidth}{>{\bfseries}l L}
    Gramática & Verbos que usam être no passé composé (movimento e mudança de estado): DR \& MRS VANDERTRAMP: Devenir, Revenir, Monter, Rester, Sortir, Venir, Aller, Naître, Descendre, Entrer, Rentrer, Tomber, Retourner, Arriver, Mourir, Partir. Concordância do participio com o sujeito: "\emph{Elle est allée.}" \\
    Vocabulário & Lista completa dos 17 verbos + seus derivados. \\
    Exercícios & 1. Conjugar verbos com être. 2. Frases com concordância. 3. Diferenciar verbos com avoir vs être. \\
    Conversação & Narrar uma viagem: "\emph{Je suis allé à Paris. Je suis arrivé lundi. Je suis resté trois jours.}" \\
    Anki & 20 cards: passé composé com être. \\
\end{tabularx}

\subsubsection*{Semana 7: Imparfait}
\begin{tabularx}{\linewidth}{>{\bfseries}l L}
    Gramática & Imparfait: radical da 1ª pessoa do plural do presente (nous parlons $\to$ parl-) + terminações: -ais, -ais, -ait, -ions, -iez, -aient. Único irregular: être (j'étais). \\
    Uso & Ações habituais no passado, descrições, cenário para outra ação. Ex: "\emph{Quand j'étais enfant, je jouais au foot.}" \\
    Exercícios & 1. Conjugar verbos no imparfait. 2. Descrever a infância. 3. Diferenciar passé composé vs imparfait. \\
    Conversação & "\emph{Quand j'étais petit, j'habitais à... Je jouais avec mes amis...}" \\
    Anki & 20 cards: imparfait. \\
\end{tabularx}

\subsubsection*{Semana 8: Passé Composé vs Imparfait (Narrativa)}
\begin{tabularx}{\linewidth}{>{\bfseries}l L}
    Gramática & Contraste: passé composé (ação concluída, pontual, sequencial) vs imparfait (ação em andamento, habitual, descritiva, cenário). Ex: "\emph{Je marchais quand j'ai vu un ami.}" \\
    Exercícios & 1. Completar histórias escolhendo o tempo correto. 2. Escrever um parágrafo misturando os dois. 3. Identificar em textos autênticos. \\
    Conversação & Contar uma história completa (ex: um dia memorável) usando ambos os tempos. \\
    Anki & 20 cards: frases com escolha de tempo. \\
\end{tabularx}

\subsection{Mês 3: Futuro, Condicional e Subjuntivo (Semanas 9–12)}
\subsubsection*{Semana 9: Futur Proche e Futur Simple}
\begin{tabularx}{\linewidth}{>{\bfseries}l L}
    Gramática & Futur proche: aller + infinitif (je vais étudier). Futur simple: infinitif + terminaisons -ai, -as, -a, -ons, -ez, -ont. Irregulares: être (ser-), avoir (aur-), aller (ir-), faire (fer-), pouvoir (pourr-), vouloir (voudr-), voir (verr-). \\
    Exercícios & 1. Formar frases com futur proche. 2. Conjugar no futur simple. 3. Planejar o futuro. \\
    Conversação & "\emph{L'année prochaine, je voyagerai en France. D'abord, je vais visiter Paris...}" \\
    Anki & 20 cards: futur. \\
\end{tabularx}

\subsubsection*{Semana 10: Conditionnel Présent}
\begin{tabularx}{\linewidth}{>{\bfseries}l L}
    Gramática & Conditionnel présent: radical do futur simple + terminaisons de l'imparfait (-ais, -ais, -ait, -ions, -iez, -aient). Usos: polidez (je voudrais), hipoteses (si j'avais...), conselhos (tu devrais). \\
    Exercícios & 1. Conjugar no conditionnel. 2. Frases de polidez. 3. Frases com "si". \\
    Conversação & Pedir educadamente: "\emph{Je voudrais un café, s'il vous plaît.}" Dar conselhos: "\emph{Tu devrais étudier plus.}" \\
    Anki & 20 cards: conditionnel. \\
\end{tabularx}

\subsubsection*{Semana 11: Subjonctif Présent (Introdução)}
\begin{tabularx}{\linewidth}{>{\bfseries}l L}
    Gramática & Formação do subjonctif présent: radical da 3ª pessoa do plural do presente (ils parlent $\to$ parl-) + terminações: -e, -es, -e, -ions, -iez, -ent. Irregulares: être (sois, sois, soit, soyons, soyez, soient); avoir (aie, aies, ait, ayons, ayez, aient); aller (aille, ailles, aille, allions, alliez, aillent); faire (fasse); pouvoir (puisse); savoir (sache). \\
    Uso & Após expressões de necessidade, vontade, emoção, dúvida: "\emph{Il faut que tu étudies. Je veux que tu viennes.}" \\
    Exercícios & 1. Conjugar no subjonctif. 2. Completar frases. \\
    Conversação & Expressar desejos e necessidades: "\emph{Il faut que je parte. Je suis content que tu sois là.}" \\
    Anki & 20 cards: subjonctif. \\
\end{tabularx}

\subsubsection*{Semana 12: Consolidação de Todos os Tempos}
\begin{tabularx}{\linewidth}{>{\bfseries}l L}
    Gramática & Revisão geral: présent, passé composé, imparfait, futur, conditionnel, subjonctif. Misturar tempos em uma conversa. \\
    Exercícios & 1. Teste escrito com todos os tempos. 2. Produzir um diálogo longo. 3. Identificar tempos em textos. \\
    Conversação & Conversa de 15 minutos sobre temas variados (viagens, planos, infância, desejos). \\
    Anki & Revisão intensiva de todos os cards. \\
\end{tabularx}

\subsection{Mês 4: Consolidação B1 (Semanas 13–16)}
\subsubsection*{Semana 13: Leitura e Vocabulário Temático}
\begin{tabularx}{\linewidth}{>{\bfseries}l L}
    Leitura & Artigos de jornal adaptados (ex: "Ijourlactu"), contos curtos (200-300 palavras). Identificar cognatos e novos vocabulários. \\
    Vocabulário & Temas: voyage, travail, santé, famille, loisirs. \\
    Exercícios & Responder perguntas de compreensão. \\
    Conversação & Discutir o texto lido. \\
    Anki & Vocabulário temático. \\
\end{tabularx}

\subsubsection*{Semana 14: Compréhension Orale (Listening)}
\begin{tabularx}{\linewidth}{>{\bfseries}l L}
    Áudio & Podcasts para aprendizes: "InnerFrench", "Français Authentique", "Coffee Break French". Músicas (Stromae, Angèle, Zaz), noticiários lentos ("Journal en français facile" da RFI). \\
    Atividades & Ouvir e anotar palavras-chave, responder perguntas, transcrever trechos. \\
    Conversação & Resumir o que ouviu. \\
    Anki & Frases ouvidas nos áudios. \\
\end{tabularx}

\subsubsection*{Semana 15: Production Écrite}
\begin{tabularx}{\linewidth}{>{\bfseries}l L}
    Escrita & Redações curtas (150-200 palavras) sobre temas: "Ma ville", "Un voyage inoubliable", "Mes projets pour l'avenir". \\
    Correção & Auto-correção ou com falante nativo (via Tandem, Hello-Talk). \\
    Conversação & Apresentar a redação oralmente. \\
    Anki & Erros corrigidos. \\
\end{tabularx}

\subsubsection*{Semana 16: Simulation DELF B1 et Révision Finale}
\begin{tabularx}{\linewidth}{>{\bfseries}l L}
    Avaliação & Prova simulada de DELF B1 (compréhension écrite, orale, production écrite et orale). \\
    Révision & Pontos fracos identificados. \\
    Célébration! & Vous avez atteint le niveau B1 en français ! \\
\end{tabularx}

% ----------------------------------------------------------------------
% Seção: Programa Estruturado de 4 Semanas (Intensivo)
% ----------------------------------------------------------------------
\section{Programa Estruturado de 4 Semanas (Intensivo)}
Se você preferir um foco mais intensivo nas primeiras 4 semanas, siga este plano:

\subsection{Semana 1: Fundações}
Foco: Saudações, être/avoir, presente básico \\
\begin{itemize}
    \item Dia 1-2: Dominar pronomes e essenciais
    \item Dia 3-4: Fundamentos de ser e ter
    \item Dia 5-6: Presente simples
    \item Dia 7: Revisão e conversações básicas
\end{itemize}

\subsection{Semana 2: Comunicação Essencial}
Foco: Verbos comuns, artigos, perguntas \\
\begin{itemize}
    \item Dia 1-2: Verbos regulares (parler, aller, faire)
    \item Dia 3-4: Artigos e negação
    \item Dia 5-6: Formação de perguntas
    \item Dia 7: Prática de diálogo
\end{itemize}

\subsection{Semana 3: Gramática e Adjetivos}
Foco: Passé composé, concordância de adjetivos \\
\begin{itemize}
    \item Dia 1-2: Passé composé
    \item Dia 3-4: Concordância de adjetivos
    \item Dia 5-6: Advérbios e frases comuns
    \item Dia 7: Escrever narrativa simples
\end{itemize}

\subsection{Semana 4: Integração e Fluência}
Foco: Futuro, frases complexas, conversação \\
\begin{itemize}
    \item Dia 1-2: Futur simple
    \item Dia 3-4: Frases preposicionais
    \item Dia 5-6: Cenários reais de conversação
    \item Dia 7: Avaliação completa
\end{itemize}

% ----------------------------------------------------------------------
% Seção: Os 100 Verbos Mais Usados
% ----------------------------------------------------------------------
\section{Os 100 Verbos Mais Usados em Francês}
Baseado em análises de frequência linguística (Corpus Lexique 3.80).

% CORREÇÃO: Agora a tabela tem 6 colunas (c l c l c l)
\begin{longtable}{c l c l c l}
    \toprule
    N° & Verbe & Traduction & N° & Verbe & Traduction \\
    \midrule
    1 & être & ser/estar & 51 & suivre & seguir \\
    2 & avoir & ter & 52 & entendre & ouvir \\
    3 & pouvoir & poder & 53 & marcher & andar/funcionar \\
    4 & faire & fazer & 54 & arriver & chegar/acontecer \\
    5 & dire & dizer & 55 & porter & levar/carregar \\
    6 & aller & ir & 56 & parler & falar \\
    7 & voir & ver & 57 & rester & ficar/permanecer \\
    8 & savoir & saber & 58 & sortir & sair \\
    9 & vouloir & querer & 59 & revenir & voltar \\
    10 & venir & vir & 60 & devenir & tornar-se \\
    11 & falloir & precisar (impessoal) & 61 & entrer & entrar \\
    12 & devoir & dever & 62 & repartir & partir novamente \\
    13 & croire & acreditar & 63 & oublier & esquecer \\
    14 & trouver & encontrar & 64 & appeler & chamar \\
    15 & donner & dar & 65 & préparer & preparar \\
    16 & prendre & pegar/tomar & 66 & jouer & jogar/tocar \\
    17 & parler & falar & 67 & lire & ler \\
    18 & aimer & amar/gostar & 68 & ouvrir & abrir \\
    19 & penser & pensar & 69 & mourir & morrer \\
    20 & regarder & olhar & 70 & compter & contar \\
    21 & comprendre & compreender & 71 & fermer & fechar \\
    22 & connaître & conhecer & 72 & perdre & perder \\
    23 & demander & perguntar/pedir & 73 & offrir & oferecer \\
    24 & répondre & responder & 74 & apprendre & aprender \\
    25 & rendre & devolver/tornar & 75 & recevoir & receber \\
    26 & attendre & esperar (aguardar) & 76 & envoyer & enviar \\
    27 & laisser & deixar (permitir) & 77 & essayer & tentar \\
    28 & passer & passar & 78 & montrer & mostrar \\
    29 & mettre & pôr/colocar & 79 & continuer & continuar \\
    30 & changer & mudar & 80 & gagner & ganhar \\
    31 & accepter & aceitar & 81 & couper & cortar \\
    32 & acheter & comprar & 82 & présenter & apresentar \\
    33 & arrêter & parar & 83 & raconter & contar (narrar) \\
    34 & apporter & trazer & 84 & quitter & deixar (partir) \\
    35 & chercher & procurar & 85 & sonner & soar/tocar \\
    36 & tomber & cair & 86 & monter & subir \\
    37 & tenir & segurar & 87 & descendre & descer \\
    38 & commencer & começar & 88 & vivre & viver \\
    39 & finir & terminar & 89 & manger & comer \\
    40 & servir & servir & 90 & boire & beber \\
    41 & oublier & esquecer & 91 & dormir & dormir \\
    42 & perdre & perder & 92 & sentir & sentir \\
    43 & rappeler & lembrar & 93 & courir & correr \\
    44 & expliquer & explicar & 94 & choisir & escolher \\
    45 & imaginer & imaginar & 95 & habiter & morar \\
    46 & espérer & esperar (ter esperança) & 96 & voyager & viajar \\
    47 & préférer & preferir & 97 & travailler & trabalhar \\
    48 & exister & existir & 98 & étudier & estudar \\
    49 & manquer & faltar & 99 & danser & dançar \\
    50 & utiliser & utilizar & 100 & chanter & cantar \\
    \bottomrule
\end{longtable}

% ----------------------------------------------------------------------
% Seção: Quick Reference
% ----------------------------------------------------------------------
\section{Os 100 Verbos Mais Usados — Quick Reference}
\subsection*{Os 20 Verbos Fundamentais - Quick Cards}
\begin{multicols}{2}
\begin{itemize}
    \item \textbf{être} - ser/estar
    \item \textbf{avoir} - ter
    \item \textbf{pouvoir} - poder
    \item \textbf{faire} - fazer
    \item \textbf{dire} - dizer
    \item \textbf{aller} - ir
    \item \textbf{voir} - ver
    \item \textbf{savoir} - saber
    \item \textbf{vouloir} - querer
    \item \textbf{venir} - vir
    \item \textbf{falloir} - precisar (imp.)
    \item \textbf{devoir} - dever
    \item \textbf{croire} - acreditar
    \item \textbf{trouver} - encontrar
    \item \textbf{donner} - dar
    \item \textbf{prendre} - pegar
    \item \textbf{parler} - falar
    \item \textbf{aimer} - amar/gostar
    \item \textbf{penser} - pensar
    \item \textbf{regarder} - olhar
\end{itemize}
\end{multicols}
\emph{Dica Estratégica:} Memorize os 20 primeiros verbos com Anki — eles cobrem aproximadamente 80\% do francês conversacional. Os 80 restantes completam habilidades específicas (+20\%).

% ----------------------------------------------------------------------
% Seção: As 100 Palavras Mais Usadas
% ----------------------------------------------------------------------
\section{As 100 Palavras Mais Usadas em Francês}
\subsection{Verbos e Palavras Essenciais (1-50)}
\begin{multicols}{2}
\begin{enumerate}[label=\arabic*., start=1]
    \item être (ser/estar)
    \item avoir (ter)
    \item aller (ir)
    \item pouvoir (poder)
    \item devoir (dever)
    \item faire (fazer)
    \item vouloir (querer)
    \item savoir (saber)
    \item prendre (pegar)
    \item venir (vir)
    \item voir (ver)
    \item donner (dar)
    \item parler (falar)
    \item mettre (pôr)
    \item trouver (encontrar)
    \item dire (dizer)
    \item croire (acreditar)
    \item arriver (chegar)
    \item appeler (chamar)
    \item chercher (procurar)
    \item le (o - m)
    \item la (a - f)
    \item de (de)
    \item un (um)
    \item une (uma)
    \item et (e)
    \item à (a)
    \item en (em)
    \item que (que)
    \item est (é)
    \item au (ao)
    \item non (não)
    \item ne (não)
    \item pas (não)
    \item se (se)
    \item qu’ (que - abreviado)
    \item son (seu - m)
    \item sa (sua - f)
    \item mes (meus)
    \item ses (seus)
    \item nous (nós)
    \item vous (vós/você)
    \item me (me)
    \item te (te)
    \item lui (lhe/ele)
    \item leur (seu/eles)
    \item les (os/as)
    \item des (dos/alguns)
    \item sur (sobre)
    \item pour (para)
\end{enumerate}
\end{multicols}

\subsection{Substantivos e Descritores Comuns (51-100)}
\begin{multicols}{2}
\begin{enumerate}[label=\arabic*., start=51]
    \item temps (tempo)
    \item façon (maneira)
    \item part (parte)
    \item chose (coisa)
    \item vie (vida)
    \item homme (homem)
    \item monde (mundo)
    \item personne (pessoa)
    \item gens (pessoas)
    \item jour (dia)
    \item main (mão)
    \item moment (momento)
    \item cas (caso)
    \item femme (mulher)
    \item lieu (lugar)
    \item idée (ideia)
    \item connaissance (conhecimento)
    \item résultat (resultado)
    \item vérité (verdade)
    \item problème (problema)
    \item raison (razão)
    \item argent (dinheiro)
    \item parole (palavra)
    \item loi (lei)
    \item nombre (número)
    \item pays (país)
    \item nom (nome)
    \item année (ano)
    \item eau (água)
    \item heure (hora)
    \item nouveau (novo)
    \item petit (pequeno)
    \item grand (grande)
    \item bon (bom)
    \item mauvais (mau)
    \item premier (primeiro)
    \item dernier (último)
    \item seul (só)
    \item possible (possível)
    \item différent (diferente)
    \item plus (mais)
    \item moins (menos)
    \item très (muito)
    \item bien (bem)
    \item mal (mal)
    \item maintenant (agora)
    \item avant (antes)
    \item aujourd’hui (hoje)
    \item demain (amanhã)
    \item hier (ontem)
\end{enumerate}
\end{multicols}

% ----------------------------------------------------------------------
% Seção: Exercícios Diários de Prática
% ----------------------------------------------------------------------
\section{Exercícios Diários de Prática}

\subsection*{Saudações Essenciais}
Instrução: Preencha os espaços com a resposta correta.
\begin{enumerate}
    \item \_\_\_ ! (Olá) $\to$ Bonjour
    \item Ça va \_\_\_ ? (Como vai você bem?) $\to$ bien
    \item \_\_\_ beaucoup (Muito obrigado) $\to$ Merci
    \item \_\_\_ -moi (Desculpe-me) $\to$ Excusez
    \item De \_\_\_ (De nada) $\to$ rien
\end{enumerate}

\subsection*{2. Conjugação de ÊTRE}
Instrução: Complete com o verbo correto.
\begin{enumerate}
    \item Je \_\_\_ français(e) $\to$ suis
    \item Tu \_\_\_ étudiant(e) $\to$ es
    \item Il \_\_\_ médecin $\to$ est
    \item Nous \_\_\_ heureux(ses) $\to$ sommes
    \item Vous \_\_\_ professeur $\to$ êtes
\end{enumerate}

\subsection{Conjunto 3: Conjugação de Avoir}
\begin{enumerate}
    \item J’ \_\_\_ un chat $\to$ ai
    \item Tu \_\_\_ une voiture $\to$ as
    \item Elle \_\_\_ des enfants $\to$ a
    \item Nous \_\_\_ une maison $\to$ avons
    \item Ils \_\_\_ des amis $\to$ ont
\end{enumerate}

\subsection{Conjunto 4: Artigos (le, la, les)}
\begin{enumerate}
    \item \_\_\_ livre $\to$ Le
    \item \_\_\_ maison $\to$ La
    \item \_\_\_ enfants $\to$ Les
    \item \_\_\_ pomme $\to$ La
    \item \_\_\_ école $\to$ L’
\end{enumerate}

\subsection{Conjunto 5: Tradução}
\begin{enumerate}
    \item I am French $\to$ Je suis français(e)
    \item She has a dog $\to$ Elle a un chien
    \item We are happy $\to$ Nous sommes heureux
    \item Do you speak English? $\to$ Parlez-vous anglais?
    \item Where is the station? $\to$ Où est la gare?
\end{enumerate}

\subsection{Drill 1: Verbos Regulares -ER}
\emph{Parler} (falar) : je parle, tu parles, il parle, nous parlons, vous parlez, ils parlent \\
Prática: Conjugue 3 verbos -ER completamente.

\subsection{Drill 2: Concordância de Adjetivos}
Masculino: un livre rouge \\
Feminino: une maison rouge \\
Masculino plural: des livres rouges \\
Feminino plural: des maisons rouges \\
Prática: Descreva 5 objetos concordando gênero/número.

\subsection{Drill 3: Negação (Ne...Pas)}
\emph{Je ne suis pas français} \\
\emph{Elle n’aime pas le café} \\
\emph{Nous ne parlons pas anglais} \\
Prática: Faça 5 frases negativas.

\subsection{Drill 4: Formação de Perguntas}
\begin{enumerate}
    \item Entonação: \emph{Tu parles français?}
    \item Est-ce que: \emph{Est-ce que tu parles français?}
    \item Inversão: \emph{Parles-tu français?}
\end{enumerate}
Prática: Crie 5 perguntas usando os 3 métodos.

\subsection{Drill 5: Contrações}
à + le = au \quad à + les = aux \quad de + le = du \quad de + les = des \\
Prática: Escreva 5 frases usando cada contração.

% ----------------------------------------------------------------------
% Seção: Vocabulário por Tema
% ----------------------------------------------------------------------
\section{Construção de Vocabulário por Tema}
\subsection{Tema 1: Comida (Nourriture)}
pain (pão), fromage (queijo), poisson (peixe), viande (carne), légumes (legumes), fruits (frutas), vin (vinho), beurre (manteiga), œufs (ovos) \\
Prática: Escreva uma refeição completa usando 8+ palavras.

\subsection{Tema 2: Família (Famille)}
père (pai), mère (mãe), frère (irmão), sœur (irmã), grand-père (avô), grand-mère (avó), oncle (tio), tante (tia) \\
Prática: Apresente sua família (5-7 frases).

\subsection{Tema 3: Casa (Pièces)}
salon (sala), chambre (quarto), cuisine (cozinha), salle de bain (banheiro), bureau (escritório), garage (garagem), jardin (jardim) \\
Prática: Descreva sua casa com localizações dos cômodos.

\subsection{Tema 4: Roupa (Vêtements)}
chemise (camisa), pantalon (calça), robe (vestido), chaussures (sapatos), chapeau (chapéu), manteau (casaco) \\
Prática: Descreva uma roupa usando cores e adjetivos.

% ----------------------------------------------------------------------
% Seção: Prática de Conversação
% ----------------------------------------------------------------------
\section{Prática de Conversação}
\subsection{Cenário 1: Encontro}
\begin{verbatim}
A: Bonjour, je m’appelle Marie.
B: Je m’appelle Pierre. Enchanté(e).
A: D’où viens-tu?
B: De Paris. Et toi?
A: De Lyon. Qu’est-ce que tu fais?
B: Je suis ingénieur. Et toi?
\end{verbatim}
Sua vez: Crie um diálogo similar.

\subsection{Cenário 2: Café}
\begin{verbatim}
Serveur: Bonjour, que désirez-vous?
Toi: Un café et un croissant, s’il vous plaît.
Serveur: Petit ou grand?
Toi: Petit, s’il vous plaît.
\end{verbatim}
Sua vez: Peça itens diferentes.

\subsection{Cenário 3: Direções}
\begin{verbatim}
Toi: Où est la gare?
Passant: À gauche, cinq minutes.
Toi: Merci beaucoup!
\end{verbatim}
Sua vez: Pergunte por diferentes locais.

% ----------------------------------------------------------------------
% Seção: Exercícios de Tradução
% ----------------------------------------------------------------------
\section{Exercícios de Tradução}
\subsection{Fácil}
\begin{enumerate}
    \item I am French $\to$ Je suis français(e)
    \item She speaks English $\to$ Elle parle anglais
    \item We have a cat $\to$ Nous avons un chat
    \item They are happy $\to$ Ils sont heureux
\end{enumerate}
\subsection{Médio}
\begin{enumerate}
    \item I don't like coffee $\to$ Je n'aime pas le café
    \item Where are you? $\to$ Où es-tu?
    \item Do they have children? $\to$ Ont-ils des enfants?
\end{enumerate}
\subsection{Avançado}
\begin{enumerate}
    \item I have spoken to him $\to$ Je lui ai parlé
    \item She was reading when I arrived $\to$ Elle lisait quand je suis arrivé(e)
    \item We went to Paris last summer $\to$ Nous sommes allés à Paris l'été dernier
\end{enumerate}

% ----------------------------------------------------------------------
% Seção: Notas Culturais
% ----------------------------------------------------------------------
\section{Notas Culturais}
\subsection{Regiões Francófonas}
\begin{itemize}
    \item França: Moda, vinho, culinária
    \item Quebec: Dialeto francês distinto
    \item Bélgica, Suíça: Sociedades multilíngues
    \item África: Senegal, Costa do Marfim, Congo
\end{itemize}
\subsection{Etiqueta}
\begin{itemize}
    \item Beijo na bochecha para cumprimentar (1-4 beijos variam por região)
    \item Sempre diga "Bonjour" e "Au revoir"
    \item Use "Vous" com estranhos/idosos
    \item Jantar é a refeição principal (19h-20h)
    \item Tempo de qualidade com a família é valorizado
\end{itemize}

% ----------------------------------------------------------------------
% Seção: Dicas de Estudo
% ----------------------------------------------------------------------
\section{Dicas de Estudo}
\subsection{Métodos}
\begin{enumerate}
    \item Shadowing de falantes nativos
    \item Parceiros de conversação
    \item Filmes e TV franceses (Netflix)
    \item Podcasts: RFI, Français Facile
    \item Escrita diária de frases
\end{enumerate}
\subsection{Truques de Memória}
\begin{enumerate}
    \item Rimar padrões verbais
    \item Agrupar verbos similares
    \item Aprender palavras em frases
    \item Visualizar significados
    \item Criar associações
\end{enumerate}
\subsection{Erros Comuns}
\begin{enumerate}
    \item Pronúncia (letras silenciosas)
    \item Esquecer o gênero dos substantivos
    \item Erros de conjugação verbal
    \item Colocação de acentos
    \item Ordem das palavras com pronomes
\end{enumerate}
\subsection{Rotina Diária}
\begin{itemize}
    \item 5 min: Vocabulário (Anki)
    \item 10 min: Drill gramatical
    \item 15 min: Prática de listening
    \item 10 min: Fala (grave-se)
    \item 5 min: Escreva 3 frases
    \item Total: 45 minutos
\end{itemize}
\subsection{Metas Semanais}
\begin{itemize}
    \item Semana 1: 25-30 palavras, dominar saudações
    \item Semana 2: Conjugações completas, perguntas
    \item Semana 3: Usar passé composé, ler textos
    \item Semana 4: Conversas de 3 minutos
\end{itemize}

% ----------------------------------------------------------------------
% Seção: Avaliação de Progresso
% ----------------------------------------------------------------------
\section{Avaliação de Progresso — Benchmarking}
\subsection{Teste Inicial - Baseline (Data: \_\_\_/\_\_\_/\_\_\_)}
\subsubsection{1. Vocabulário Básico (A1)}
Traduza sem olhar:
\begin{enumerate}
    \item Chat = \_\_\_\_\_\_\_\_\_\_
    \item Maison = \_\_\_\_\_\_\_\_\_\_
    \item Manger = \_\_\_\_\_\_\_\_\_\_
    \item Eau = \_\_\_\_\_\_\_\_\_\_
    \item Merci = \_\_\_\_\_\_\_\_\_\_
\end{enumerate}
Pontuação: \_\_\_ / 5

\subsubsection{2. Compreensão (Listening A1)}
Quando você compreende:
\begin{itemize}
    \item Saudações (Bonjour, Ça va?)
    \item Números (0-10)
    \item Nomes
    \item Comida comum
    \item Funções (Oui, Non, S’il vous plaît)
\end{itemize}
Nível Atingido: \_\_\_ / 5

\subsubsection{3. Fala Básica}
Você tem capacidade de dizer (em voz alta, confiante):
\begin{itemize}
    \item "Je m’appelle..."
    \item "Je suis de..."
    \item "J’ai... ans"
    \item "Comment allez-vous?"
    \item "D’où venez-vous?"
\end{itemize}
Nível Atingido: \_\_\_ / 5

\subsubsection{4. Conversação Simples}
Complete sem ajuda:
\begin{verbatim}
A: "Bonjour, comment t’appelles-tu?"
B: "_________________________________"
A: "D’où es-tu?"
B: "_________________________________"
A: "Que fais-tu?"
B: "_________________________________"
\end{verbatim}
Pontuação: \_\_\_ / 3

\subsection{Progressão Mensal}
\subsubsection{Mês 1 (Data: \_\_\_/\_\_\_/\_\_\_)}
\begin{tabular}{lcc}
    \toprule
    Teste & Pontuação & Nível \\
    \midrule
    Vocabulário (50 palavras) & \_\_\_ / 50 & \\
    Listening (A1) & \_\_\_ / 10 & \\
    Fala (5 frases) & \_\_\_ / 5 & \\
    Conversação (3 Q\&R) & \_\_\_ / 3 & \\
    TOTAL & \_\_\_ / 68 & \\
    \bottomrule
\end{tabular}

\subsubsection{Mês 2 (Data: \_\_\_/\_\_\_/\_\_\_)}
\begin{tabular}{lcc}
    \toprule
    Teste & Pontuação & Nível \\
    \midrule
    Vocabulário (100 palavras) & \_\_\_ / 100 & \\
    Listening (A2) & \_\_\_ / 15 & \\
    Fala (8 frases) & \_\_\_ / 8 & \\
    Conversação (5 Q\&R) & \_\_\_ / 5 & \\
    TOTAL & \_\_\_ / 128 & \\
    \bottomrule
\end{tabular}

\subsubsection{Mês 3 (Data: \_\_\_/\_\_\_/\_\_\_)}
\begin{tabular}{lcc}
    \toprule
    Teste & Pontuação & Nível \\
    \midrule
    Vocabulário (150 palavras) & \_\_\_ / 150 & \\
    Listening (B1) & \_\_\_ / 20 & \\
    Fala (12 frases) & \_\_\_ / 12 & \\
    Conversação (8 Q\&R) & \_\_\_ / 8 & \\
    TOTAL & \_\_\_ / 190 & \\
    \bottomrule
\end{tabular}

\subsection{Indicadores de Progresso}
\subsubsection*{Sinais Positivos:}
\begin{itemize}
    \item Reconheço 30+ palavras novas por mês
    \item Compreendo conversas simples em filmes
    \item Posso me apresentar sem estresse
    \item Peço no restaurante sem medo
    \item Minha pronúncia melhora a cada semana
\end{itemize}
\subsubsection*{Sinais de Alerta:}
\begin{itemize}
    \item Nenhuma palavra nova aprendida
    \item Ainda não consigo dizer meu nome em francês
    \item Tenho medo de falar em voz alta
    \item Assisto filme mas não entendo NADA
    \item Estudo mas nunca pratico
\end{itemize}

% ----------------------------------------------------------------------
% Seção: Cognatos e Falsos Amigos
% ----------------------------------------------------------------------
\section{Listas de Palavras Similares (Cognatos)}
\subsection{Cognatos Perfeitos (Escrita Idêntica ou Quase Idêntica)}
\begin{multicols}{4}
\begin{itemize}
    \item animal
    \item terrible
    \item note
    \item important
    \item flexible
    \item texte
    \item possible
    \item responsable
    \item total
    \item problème
    \item art
    \item triple
    \item général
    \item part
    \item universel
    \item naturel
    \item classe
    \item vertical
    \item original
    \item phrase
    \item culture
    \item particulier
    \item image
    \item nature
    \item visible
    \item ligne
    \item aventure
    \item horrible
    \item musique
    \item lecture
\end{itemize}
\end{multicols}

\subsection{Padrões de Conversão Mais Comuns}
\subsubsection{Português -ção $\to$ Français -tion}
\begin{tabular}{ll}
    \toprule
    Português & Français \\
    \midrule
    nação & nation \\
    canção & chanson (exceção) \\
    educação & éducation \\
    informação & information \\
    situação & situation \\
    atenção & attention \\
    condição & condition \\
    tradição & tradition \\
    produção & production \\
    construção & construction \\
    \bottomrule
\end{tabular}

\subsubsection{Português -dade $\to$ Français -té}
\begin{tabular}{ll}
    \toprule
    Português & Français \\
    \midrule
    cidade & cité \\
    verdade & vérité \\
    universidade & université \\
    felicidade & félicité (menos usado) / bonheur \\
    liberdade & liberté \\
    oportunidade & opportunité \\
    facilidade & facilité \\
    dificuldade & difficulté \\
    amizade & amitié \\
    \bottomrule
\end{tabular}

\subsubsection{Português -vel $\to$ Français -ble}
\begin{tabular}{ll}
    \toprule
    Português & Français \\
    \midrule
    possível & possible \\
    amável & aimable \\
    terrível & terrible \\
    horrível & horrible \\
    flexível & flexible \\
    responsável & responsable \\
    notável & notable \\
    durável & durable \\
    favorável & favorable \\
    viável & viable \\
    \bottomrule
\end{tabular}

\subsection{Falsos Amigos (Faux Amis) — Cuidado!}
\begin{tabular}{lll}
    \toprule
    Português & Francês & Significado real em francês \\
    \midrule
    atualmente & actuellement & atualmente (não "atualmente") \\
    biblioteca & bibliothèque & biblioteca \\
    presentear & présenter & apresentar (não dar presente) \\
    propaganda & propagande & propaganda política (não publicidade) \\
    saco & sac & bolsa/saco \\
    puxe & pousser & empurrar (não puxar) \\
    tire & tirer & puxar (não tirar) \\
    brincar & brinquebaler & sacudir (não brincar) \\
    barato & baratin & confusão (não barato) \\
    esquisito & exquis & delicioso, requintado \\
    rato & rate & falha (não rato) \\
    vaso & vase & vaso (mas também "lama") \\
    assistir & assister & participar (não assistir como espectador) \\
    \bottomrule
\end{tabular}

% ----------------------------------------------------------------------
% Seção: Recursos Recomendados
% ----------------------------------------------------------------------
\section{Recursos Recomendados}
\begin{itemize}
    \item \textbf{Dicionários:} WordReference (com fórum), Linguee, Reverso Context
    \item \textbf{Anki Decks:} "French Frequency 500", "French Verb Conjugations", "DELF B1 Preparation"
    \item \textbf{YouTube:} Lingoni French, Easy French, Français avec Pierre, InnerFrench
    \item \textbf{Podcasts:} "InnerFrench", "Français Authentique", "Coffee Break French", "Journal en français facile" (RFI)
    \item \textbf{Séries e Filmes:} "Extra French" (A2), "Dix pour cent" (B1+), "Lupin" (B1+)
    \item \textbf{Apps de Troca de Idiomas:} Tandem, HelloTalk
\end{itemize}

% ----------------------------------------------------------------------
% Seção: Conselhos Finais
% ----------------------------------------------------------------------
\section{Conselhos Finais}
\begin{enumerate}
    \item A regularidade é mais importante que a intensidade. 30 minutos por dia valem mais que 3 horas uma vez por semana.
    \item Nunca pule o Anki. É a base da memorização de longo prazo.
    \item Fale em voz alta desde o primeiro dia. Mesmo que cometa erros.
    \item Ouça muito francês. Seu cérebro precisa se habituar aos sons.
    \item Seja paciente. O francês é mais desafiador para um lusófono, mas com este plano, você atingirá o B1 em 4 meses.
\end{enumerate}

\end{document}